% ============================================================================
% WORKBOOK — Section 1.1: Unit Rates with Fractions
% CCSS 7.RP.A.1
% Practice-focused reinforcement for the study guide topic
% ============================================================================
\section{Unit Rates with Fractions}
\topicTitle{Unit Rates with Fractions}

\begin{quickReview}{Unit Rates with Fractions}
A \textbf{unit rate} tells you how much of one quantity per \textbf{one unit} of another.

\begin{itemize}[leftmargin=*]
    \item To find a unit rate, \textbf{divide}: $\text{unit rate} = \dfrac{\text{first quantity}}{\text{second quantity}}$.
    \item When quantities are fractions, use \textbf{Keep--Change--Flip}: $\dfrac{a}{b} \div \dfrac{c}{d} = \dfrac{a}{b} \times \dfrac{d}{c}$.
\end{itemize}

\medskip
\textbf{Example:} A runner covers $\frac{3}{4}$ mile in $\frac{1}{2}$ hour. Unit rate $= \frac{3}{4} \div \frac{1}{2} = \frac{3}{4} \times \frac{2}{1} = \frac{3}{2} = 1\frac{1}{2}$ miles per hour.
\end{quickReview}

% ── Warm-Up ──────────────────────────────────────────────────────────────────
\begin{practiceBox}[funGreen]{\faSun~Warm-Up}

\practiceHeader[funGreen]{Find the Unit Rate (whole numbers)}

\begin{multicols}{2}
\prob $120$ miles in $4$ hours \answerBlank[3cm]
\answerExplain{$30$ miles per hour}{To find the unit rate, divide the total distance by the total time: $120 \div 4 = 30$ miles per hour.}

\prob $\$45$ for $5$ shirts \answerBlank[3cm]
\answerExplain{$\$9$ per shirt}{Divide total cost by the number of shirts: $45 \div 5 = 9$. Each shirt costs $\$9$.}

\prob $84$ pages in $7$ days \answerBlank[3cm]
\answerExplain{$12$ pages per day}{Divide total pages by the number of days: $84 \div 7 = 12$ pages per day.}

\prob $60$ apples in $12$ bags \answerBlank[3cm]
\answerExplain{$5$ apples per bag}{Divide the total by the number of groups: $60 \div 12 = 5$ apples per bag.}

\prob $\$32$ for $8$ gallons \answerBlank[3cm]
\answerExplain{$\$4$ per gallon}{Divide total cost by gallons: $32 \div 8 = 4$. Each gallon costs $\$4$.}

\prob $200$ words in $5$ minutes \answerBlank[3cm]
\answerExplain{$40$ words per minute}{Divide total words by minutes: $200 \div 5 = 40$ words per minute.}
\end{multicols}
\end{practiceBox}

% ── Practice A: Fraction Unit Rates ──────────────────────────────────────────
\begin{practiceBox}{Computing Unit Rates with Fractions}

\practiceHeader{Divide to find each unit rate. Simplify your answer.}

\prob A baker uses $\frac{2}{3}$ cup of flour for $\frac{1}{3}$ batch. How much flour per batch?
\answerBlank[3cm]
\answerExplain{$2$ cups per batch}{Divide cups by batches: $\frac{2}{3} \div \frac{1}{3} = \frac{2}{3} \times \frac{3}{1} = \frac{6}{3} = 2$ cups per batch.}

\prob A hose fills $\frac{3}{4}$ of a bucket in $\frac{1}{2}$ minute. How many buckets per minute?
\answerBlank[3cm]
\answerExplain{$1\frac{1}{2}$ buckets per minute}{Divide buckets by minutes: $\frac{3}{4} \div \frac{1}{2} = \frac{3}{4} \times \frac{2}{1} = \frac{6}{4} = \frac{3}{2} = 1\frac{1}{2}$ buckets per minute.}

\prob A train covers $\frac{5}{6}$ mile in $\frac{2}{3}$ minute. What is the speed in miles per minute?
\answerBlank[3cm]
\answerExplain{$1\frac{1}{4}$ miles per minute}{Divide distance by time: $\frac{5}{6} \div \frac{2}{3} = \frac{5}{6} \times \frac{3}{2} = \frac{15}{12} = \frac{5}{4} = 1\frac{1}{4}$ miles per minute.}

\prob A machine packages $\frac{4}{5}$ kg of rice in $\frac{2}{5}$ hour. How many kg per hour?
\answerBlank[3cm]
\answerExplain{$2$ kg per hour}{Divide kg by hours: $\frac{4}{5} \div \frac{2}{5} = \frac{4}{5} \times \frac{5}{2} = \frac{20}{10} = 2$ kg per hour.}

\prob A gardener plants $\frac{7}{8}$ row in $\frac{3}{4}$ hour. How many rows per hour?
\answerBlank[3cm]
\answerExplain{$1\frac{1}{6}$ rows per hour}{Divide rows by hours: $\frac{7}{8} \div \frac{3}{4} = \frac{7}{8} \times \frac{4}{3} = \frac{28}{24} = \frac{7}{6} = 1\frac{1}{6}$ rows per hour.}

\prob A recipe calls for $\frac{1}{3}$ cup of sugar for $\frac{5}{6}$ batch. How much sugar per batch?
\answerBlank[3cm]
\answerExplain{$\frac{2}{5}$ cup per batch}{Divide cups by batches: $\frac{1}{3} \div \frac{5}{6} = \frac{1}{3} \times \frac{6}{5} = \frac{6}{15} = \frac{2}{5}$ cup per batch.}

\end{practiceBox}

% ── Practice B: Comparing Rates ──────────────────────────────────────────────
\begin{practiceBox}[funTeal]{Who Is Faster? Which Costs Less?}

\practiceHeader[funTeal]{Compare by finding the unit rate for each. Circle your answer.}

\prob Runner A covers $\frac{3}{5}$ mile in $\frac{1}{4}$ hour. Runner B covers $\frac{5}{6}$ mile in $\frac{1}{3}$ hour. Who is faster?
\answerBlank[3cm]
\answerExplain{Runner B ($2\frac{1}{2}$ mph $>$ $2\frac{2}{5}$ mph)}{Runner A: $\frac{3}{5} \div \frac{1}{4} = \frac{3}{5} \times 4 = \frac{12}{5} = 2\frac{2}{5}$ mph. Runner B: $\frac{5}{6} \div \frac{1}{3} = \frac{5}{6} \times 3 = \frac{15}{6} = 2\frac{1}{2}$ mph. Since $2\frac{1}{2} = 2.5 > 2\frac{2}{5} = 2.4$, Runner B is faster.}

\prob Brand X: $\frac{3}{4}$ lb of cheese for $\$\frac{5}{2}$. Brand Y: $\frac{2}{3}$ lb for $\$2$. Which costs less per pound?
\answerBlank[3cm]
\answerExplain{Brand Y ($\$3$/lb $<$ $\$3\frac{1}{3}$/lb)}{Brand X: $\frac{5}{2} \div \frac{3}{4} = \frac{5}{2} \times \frac{4}{3} = \frac{20}{6} = \$3\frac{1}{3}$ per lb. Brand Y: $2 \div \frac{2}{3} = 2 \times \frac{3}{2} = \$3$ per lb. Brand Y is cheaper.}

\prob Pipe A fills $\frac{1}{2}$ tank in $\frac{2}{5}$ hour. Pipe B fills $\frac{3}{4}$ tank in $\frac{1}{2}$ hour. Which pipe fills faster?
\answerBlank[3cm]
\answerExplain{Pipe B ($1\frac{1}{2}$ tanks/hr $>$ $1\frac{1}{4}$ tanks/hr)}{Pipe A: $\frac{1}{2} \div \frac{2}{5} = \frac{1}{2} \times \frac{5}{2} = \frac{5}{4} = 1\frac{1}{4}$ tanks/hr. Pipe B: $\frac{3}{4} \div \frac{1}{2} = \frac{3}{4} \times 2 = \frac{3}{2} = 1\frac{1}{2}$ tanks/hr. Pipe B fills faster.}

\prob Cyclist A travels $\frac{7}{10}$ mile in $\frac{1}{5}$ hour. Cyclist B travels $\frac{4}{5}$ mile in $\frac{1}{4}$ hour. Who has a greater speed?
\answerBlank[3cm]
\answerExplain{Cyclist A ($3\frac{1}{2}$ mph $>$ $3\frac{1}{5}$ mph)}{Cyclist A: $\frac{7}{10} \div \frac{1}{5} = \frac{7}{10} \times 5 = \frac{35}{10} = 3\frac{1}{2}$ mph. Cyclist B: $\frac{4}{5} \div \frac{1}{4} = \frac{4}{5} \times 4 = \frac{16}{5} = 3\frac{1}{5}$ mph. Cyclist A is faster.}

\end{practiceBox}

% ── Practice C: True or False ────────────────────────────────────────────────
\begin{practiceBox}[funPurple]{True or False?}

\trueOrFalse{To find a unit rate with fractions, you multiply the two fractions.}
\answerExplain{False}{To find a unit rate, you \emph{divide} the first quantity by the second. Dividing fractions uses Keep--Change--Flip (multiply by the reciprocal), but the operation is division, not just multiplication.}

\trueOrFalse{$\frac{2}{3} \div \frac{1}{6} = 4$.}
\answerExplain{True}{Keep--Change--Flip: $\frac{2}{3} \times \frac{6}{1} = \frac{12}{3} = 4$. The statement is correct.}

\trueOrFalse{If you walk $\frac{1}{2}$ mile in $\frac{1}{4}$ hour, your unit rate is $\frac{1}{8}$ mile per hour.}
\answerExplain{False}{Divide distance by time: $\frac{1}{2} \div \frac{1}{4} = \frac{1}{2} \times 4 = 2$ miles per hour, not $\frac{1}{8}$. Multiplying $\frac{1}{2} \times \frac{1}{4} = \frac{1}{8}$ is the wrong operation.}

\end{practiceBox}

% ── Practice D: Word Problems ────────────────────────────────────────────────
\begin{practiceBox}[funBlue]{\faGlobe~Word Problems}

\wordProblem{A painter can paint $\frac{2}{5}$ of a fence in $\frac{3}{4}$ hour. At this rate, how much of the fence can the painter finish in one hour?}{of the fence}
\answerExplain{$\frac{8}{15}$ of the fence}{Find the unit rate: $\frac{2}{5} \div \frac{3}{4} = \frac{2}{5} \times \frac{4}{3} = \frac{8}{15}$ of the fence per hour.}

\wordProblem{A juice machine fills $\frac{5}{6}$ of a bottle every $\frac{2}{3}$ second. How many bottles does it fill per second?}{bottles per second}
\answerExplain{$1\frac{1}{4}$ bottles per second}{Divide bottles by seconds: $\frac{5}{6} \div \frac{2}{3} = \frac{5}{6} \times \frac{3}{2} = \frac{15}{12} = \frac{5}{4} = 1\frac{1}{4}$ bottles per second.}

\wordProblem{Sam reads $\frac{3}{4}$ of a chapter in $\frac{1}{2}$ hour. Emma reads $\frac{2}{3}$ of a chapter in $\frac{1}{3}$ hour. Who reads faster, and by how much?}{chapters per hour}
\answerExplain{Emma is faster by $\frac{1}{2}$ chapter per hour}{Sam: $\frac{3}{4} \div \frac{1}{2} = \frac{3}{4} \times 2 = \frac{3}{2} = 1\frac{1}{2}$ ch/hr. Emma: $\frac{2}{3} \div \frac{1}{3} = \frac{2}{3} \times 3 = 2$ ch/hr. Difference: $2 - 1\frac{1}{2} = \frac{1}{2}$ chapter per hour.}

\end{practiceBox}

% ── Challenge ────────────────────────────────────────────────────────────────
\begin{challengeBox}
\prob A snail travels $\frac{3}{8}$ meter in $\frac{5}{6}$ hour. A turtle travels $\frac{7}{12}$ meter in $1\frac{1}{4}$ hours. Which animal has a greater speed? By how much?
\answerExplain{Snail: $\frac{9}{20}$ m/hr $=$ Turtle: $\frac{7}{15}$ m/hr; Turtle is faster by $\frac{1}{60}$ m/hr}{Snail: $\frac{3}{8} \div \frac{5}{6} = \frac{3}{8} \times \frac{6}{5} = \frac{18}{40} = \frac{9}{20}$ m/hr. Turtle: $\frac{7}{12} \div \frac{5}{4} = \frac{7}{12} \times \frac{4}{5} = \frac{28}{60} = \frac{7}{15}$ m/hr. Compare: $\frac{9}{20} = \frac{27}{60}$ and $\frac{7}{15} = \frac{28}{60}$. The turtle is faster by $\frac{1}{60}$ m/hr.}

\prob A factory machine A produces $\frac{5}{6}$ widget in $\frac{3}{8}$ hour. Machine B produces $1\frac{1}{3}$ widgets in $\frac{2}{3}$ hour. Which machine is more productive?
\answerBlank[3cm]
\answerExplain{Machine A ($2\frac{2}{9}$ widgets/hr $>$ $2$ widgets/hr)}{Machine A: $\frac{5}{6} \div \frac{3}{8} = \frac{5}{6} \times \frac{8}{3} = \frac{40}{18} = \frac{20}{9} = 2\frac{2}{9}$ widgets/hr. Machine B: $\frac{4}{3} \div \frac{2}{3} = \frac{4}{3} \times \frac{3}{2} = \frac{12}{6} = 2$ widgets/hr. Machine A is more productive.}
\end{challengeBox}

\encouragement{Awesome work with fraction unit rates! You are becoming a unit rate expert!}
